\documentclass{stex}

\libusepackage{naproche}
\libusepackage{copyright}
\libinput{preamble}

\title{Furstenberg's Proof of the Infinitude of Primess in \Naproche}
\author{Andrei Paskevich, Marcel Schütz}
\date{2025}

\begin{document}
\begin{smodule}{furstenberg.ftl}

\maketitle

\importmodule[libraries/arithmetics]{primes.ftl}
\importmodule[libraries/arithmetics]{nat-is-a-set.ftl}
\importmodule[libraries/set-theory]{zfc.ftl}
\importmodule[libraries/foundations]{closure-under-finite-unions.ftl}

\symdef{Compl}[args=1]{#1^{\comp\complement}}
\symdef{ArithSeq}[args=2]{\comp{\mathrm{N}}_{#1\comp,#2}}
\symdef{Int}{\mathbb{Z}}
\symdef{PrimesSeq}{\mathrm{P}}

\symdecl*{open}
\symdecl*{system of open sets}
\symdecl*{closed}
\symdecl*{system of closed sets}
\symdecl*{Furstenberg}

\noindent This is a formalization of Furstenberg's topological proof of the
infinitude of primes \cite[p. 353]{Furstenberg1955}.

The central idea of Furstenberg's proof is to define a certain topology on
$\Naturals$ from the properties of which we can deduce that the set of
primes is infinite.\footnote{Actually, Furstenberg's proof makes use of a
topology on $\Int$. But this topology can as well be restricted to
$\Naturals$ without substantially changing the proof.}

\begin{forthel}
  Let $n, m, k$ denote natural numbers.
  Let $p, q$ denote nonzero natural numbers.

  \begin{definition}[for=Compl]
    Let $A$ be a subset of $\Naturals$.
    $\Compl{A} \DefEq \Naturals \Diff A$.
  \end{definition}

  Let the complement of $A$ stand for $\Compl{A}$.

  \begin{lemma}
    The complement of any subset of $\Naturals$ is a subset of $\Naturals$.
  \end{lemma}
\end{forthel}

Towards a suitable topology on $\Naturals$ let us define \textit{arithmetic
sequences} $\ArithSeq{n}{q}$ on $\Naturals$.

\begin{forthel}
  \begin{definition}[for=ArithSeq]
    $\ArithSeq{n}{q} \DefEq \SClass{m}{\Naturals}{\Cong{m}{n}{q}}$.
  \end{definition}
\end{forthel}

This allows us to define the \textit{evenly spaced natural number
topology} on $\Naturals$, whose open sets are defined as follows.

\begin{forthel}
  \begin{definition}[for=open]
    Let $U$ be a subset of $\Naturals$.
    $U$ is open iff for any $n \Elem U$ there exists a $q$ such that
    $\ArithSeq{n}{q} \Subclass U$.
  \end{definition}

  \begin{definition}[for=system of open sets]
    A system of open sets is a system of sets $S$ such that every element of
    $S$ is an open subset of $\Naturals$.
  \end{definition}

  \begin{lemma}
    Every system of open sets is a set.
  \end{lemma}
  \begin{proof}
    Let $S$ be a system of open sets.
    Then $S \Subclass \Pow(\Naturals)$.
    Hence $S$ is a set.
  \end{proof}
\end{forthel}

We can show that the open sets indeed form a topology on $\Naturals$.

\begin{forthel}
  \begin{lemma}
    $\Naturals$ and $\EmptyClass$ are open.
  \end{lemma}

  \begin{lemma}
    Let $U,V$ be open subsets of $\Naturals$.
    Then $U \Intersect V$ is open.
  \end{lemma}
  \begin{proof}
    Let $n \Elem U \Intersect V$.
    Take a $q$ such that $\ArithSeq{n}{q} \Subclass U$.
    Take a $p$ such that $\ArithSeq{n}{p} \Subclass V$.
    Then $p \Mul q \NotEq \Zero$.

    Let us show that $\ArithSeq{n}{p \Mul q} \Subclass U \Intersect V$.
      Let $m \Elem \ArithSeq{n}{p \Mul q}$.
      We have\linebreak $\Cong{m}{n}{p \Mul q}$.
      Hence $\Cong{m}{n}{p}$ and $\Cong{m}{n}{q}$.
      Thus $m \Elem \ArithSeq{n}{p}$ and $m \Elem \ArithSeq{n}{q}$.
      Therefore $m \Elem U$ and $m \Elem V$.
      Consequently $m \Elem U \Intersect V$.
    End.
  \end{proof}

  \begin{lemma}
    Let $S$ be a system of open sets.
    Then $\UnionOver S$ is open.
  \end{lemma}
  \begin{proof}
    Let $n \Elem \UnionOver S$.
    Take a set $M$ such that $n \Elem M \Elem S$.
    Consider a $q$ such that $\ArithSeq{n}{q} \Subclass M$.
    Then $\ArithSeq{n}{q} \Subclass \UnionOver S$.
  \end{proof}
\end{forthel}

Now that we have a topology of open sets on $\Naturals$, we can continue
with a characterization of closed sets whose key property is that they are
closed under finite unions.

\begin{forthel}
  \begin{definition}[for=closed]
    Let $A$ be a subset of $\Naturals$.
    $A$ is closed iff $\Compl{A}$ is open.
  \end{definition}

  \begin{definition}[for=system of closed sets]
    A system of closed sets is a system of sets $S$ such that every element of
    $S$ is a closed subset of $\Naturals$.
  \end{definition}

  \begin{lemma}
    Every system of closed sets is a set.
  \end{lemma}
  \begin{proof}
    Let $S$ be a system of closed sets.
    Then $S \Subclass \Pow(\Naturals)$.
    $\Pow(\Naturals)$ is a set.
    Hence $S$ is a set.
  \end{proof}

  \begin{lemma}
    Let $S$ be a finite system of closed sets.
    Then $\UnionOver S$ is closed.
  \end{lemma}
  \begin{proof}
    Define $C \DefEq \CClass{X}{X \text{ is a closed subset of } \Naturals}$.

    Let us show that $A \Union B \Elem C$ for any $A, B \Elem C$.
      Let $A, B \Elem C$.
      Then $A, B$ are closed subsets of $\Naturals$.
      We have $(\Compl{A \Union B}) \Eq \Compl{A} \Intersect \Compl{B}$. %!
      $\Compl{A}$ and $\Compl{B}$ are open.
      Hence $\Compl{A} \Intersect \Compl{B}$ is open.
      Thus $A \Union B$ is a closed subset of $\Naturals$.
    End.

    Therefore $C$ is closed under finite unions.
    Consequently $\UnionOver S \Elem C$.
    Indeed $S$ is a subset of $C$.
  \end{proof}
\end{forthel}

An important step towards Furstenberg's proof is to show that arithmetic
sequences are closed.

\begin{forthel}
  \begin{lemma}
    $\ArithSeq{n}{q}$ is closed.
  \end{lemma}
  \begin{proof}
    Let $m \Elem \Compl{\ArithSeq{n}{q}}$.

    Let us show that $\ArithSeq{m}{q} \Subclass \Compl{\ArithSeq{n}{q}}$.
      Let $k \Elem \ArithSeq{m}{q}$.
      Assume $k \NotElem \Compl{\ArithSeq{n}{q}}$.
      Then $\Cong{k}{m}{q}$ and $\Cong{n}{k}{q}$.
      Hence $\Cong{m}{n}{q}$.
      Therefore $m \Elem \ArithSeq{n}{q}$.
      Contradiction.
    End.
  \end{proof}
\end{forthel}

Identifying each prime number $p$ with the arithmetic sequence $\ArithSeq{0}{p}$
yields a bijection between the set $\Primes$ of all prime numbers and the set
$\PrimesSeq$ of all such sequences $\ArithSeq{0}{p}$.
Thus to show that there are infinitely many primes it suffices to show that
$\PrimesSeq$ is infinite.

\begin{forthel}
  \begin{definition}[for=PrimesSeq]
    $\PrimesSeq \DefEq \CClass{\ArithSeq{\Zero}{p}}{p \Elem \Primes}$.
  \end{definition}

  \begin{lemma}
    $\PrimesSeq$ is a system of closed sets.
  \end{lemma}
  \begin{proof}
    $\ArithSeq{\Zero}{p}$ is a closed subset of $\Naturals$ for every $p \Elem \Primes$.
  \end{proof}

  \begin{lemma}
    $\PrimesSeq$ is a set that is equinumerous to $\Primes$.
  \end{lemma}
  \begin{proof}
    (1) $\PrimesSeq$ is a set.
    Indeed $\PrimesSeq \Subclass \Pow(\Naturals)$.

    (2) $\PrimesSeq$ is equinumerous to $\Primes$.
    \begin{proof}
      Define $f(p) \DefEq \ArithSeq{\Zero}{p}$ for $p \Elem \Primes$.

      Let us show that $f$ is injective.
        Let $p, q \Elem \Primes$.
        Assume $f(p) \Eq f(q)$.
        Then $\ArithSeq{\Zero}{p} \Eq \ArithSeq{\Zero}{q}$.
        We have $\ArithSeq{\Zero}{p} \Eq \SClass{m}{\Naturals}{\Cong{m}{\Zero}{p}}$ and
        $\ArithSeq{\Zero}{q} \Eq\linebreak \SClass{m}{\Naturals}{\Cong{m}{\Zero}{q}}$.

        We can show that for all $m \Elem \Naturals$ we have $p \Divides m$ iff $q \Divides m$.
          Let $m \Elem \Naturals$.
          Then $\Cong{m}{\Zero}{p}$ iff $\Cong{m}{\Zero}{q}$.
          Thus $m \Mod p \Eq \Zero \Mod p$ iff $m \Mod q \Eq\linebreak \Zero \Mod q$.
          We have $\Zero \Mod p \Eq \Zero \Eq \Zero \Mod q$.
          Hence $m \Mod p \Eq \Zero$ iff $m \Mod q \Eq\linebreak \Zero$.
          Therefore $p \Divides m$ iff $q \Divides m$.
        End.

        Consequently $p \Eq q$.
      End.

      $f$ is surjective onto $\PrimesSeq$.
      Thus $f$ is a bijection between $\Primes$ and $\PrimesSeq$.
    \end{proof}
  \end{proof}

  \begin{theorem}[title=Furstenberg,name=Furstenberg]
    $\Primes$ is infinite.
  \end{theorem}
  \begin{proof}
    $\UnionOver \PrimesSeq$ is a subset of $\Naturals$.

    Let us show that for any $n \Elem \Naturals$ we have $n \Elem \UnionOver \PrimesSeq$ iff $n$
    has a prime divisor.
      Let $n \Elem \Naturals$.

      If $n$ has a prime divisor then $n$ belongs to $\UnionOver \PrimesSeq$.
      \begin{proof}
        Assume $n$ has a prime divisor.
        Take a prime divisor $p$ of $n$.
        We have $\ArithSeq{\Zero}{p} \Elem \PrimesSeq$.
        Hence $n \Elem \ArithSeq{\Zero}{p}$.
      \end{proof}

      If $n$ belongs to $\UnionOver \PrimesSeq$ then $n$ has a prime divisor.
      \begin{proof}
        Assume that $n$ belongs to $\UnionOver \PrimesSeq$.
        Take a prime number $r$ such that $n \Elem \ArithSeq{\Zero}{r}$.
        Hence $\Cong{n}{\Zero}{r}$.
        Thus $n \Mod r \Eq \Zero \Mod r \Eq \Zero$.
        Therefore $r$ is a prime divisor of $n$.
      \end{proof}
    End.

    Hence For all $n \Elem \Naturals$ we have $n \Elem \Compl{\UnionOver \PrimesSeq}$ iff
    $n$ has no prime divisor.
    $\One$ has no prime divisor and any natural number having no prime
    divisor is equal to $\One$.
    Therefore $\Compl{\UnionOver \PrimesSeq} \Eq \Singleton{\One}$.
    Indeed $\Compl{\UnionOver \PrimesSeq} \Subclass \Singleton{\One}$ and $\Singleton{\One}
    \Subclass \Compl{\UnionOver \PrimesSeq}$.

    $\PrimesSeq$ is infinite.
    \begin{proof}[method=contradiction]
      Assume that $\PrimesSeq$ is finite.
      Then $\UnionOver \PrimesSeq$ is closed and $\Compl{\UnionOver \PrimesSeq}$ is open.
      Take a $p$ such that $\ArithSeq{\One}{p} \Subclass \Compl{\UnionOver \PrimesSeq}$.
      $\One \Plus p$ is an element of $\ArithSeq{\One}{p}$.
      Indeed $\Cong{\One \Plus p}{\One}{p}$
      (by \sn{congruency of sum}).
      $\One \Plus p$ is not equal to $\One$.
      Hence $\One \Plus p \NotElem \Compl{\UnionOver \PrimesSeq}$.
      Contradiction.
    \end{proof}
  \end{proof}
\end{forthel}

\printbibliography
\printcopyright{Andrei Paskevich (2008), Marcel Schütz (2021--2025)}

\end{smodule}
\end{document}
